\documentclass{article}
\usepackage{amsmath,amssymb,amsthm}
\usepackage{algorithm}
\usepackage{algpseudocode}
\usepackage{geometry}
\geometry{margin=1in}
\newtheorem{theorem}{Theorem}
\newtheorem{lemma}{Lemma}
\newtheorem{assumption}{Assumption}
\newcommand{\E}{\mathbb{E}}
\newcommand{\R}{\mathbb{R}}
\newcommand{\norm}[1]{\left\|#1\right\|}
\newcommand{\inner}[1]{\left\langle#1\right\rangle}
\newcommand{\prox}{\operatorname{prox}}
\begin{document}

%%%%%%%%%%%%%%%%%%%%%%%%%%%%%%%%%%%%%%%%%%%%%%%%%%%%%%%%%%%%
\section*{AdamW (Adam with Decoupled Weight Decay)}
%%%%%%%%%%%%%%%%%%%%%%%%%%%%%%%%%%%%%%%%%%%%%%%%%%%%%%%%%%%%

%%%%%%%%%%%%%%%%%%%%%%%%%%%%%%%%%%%%%%%%%%%%%%%%%%%%%%%%%%%%
\section*{Mathematical Formulation}
%%%%%%%%%%%%%%%%%%%%%%%%%%%%%%%%%%%%%%%%%%%%%%%%%%%%%%%%%%%%
Let $f:\R^{d}\rightarrow\R$ be the objective function.  
At iteration $t\;(t=0,1,2,\dots)$ the algorithm stores:

\begin{itemize}
  \item First‑moment estimate 
        \[
        m_{t}= \beta_{1}m_{t-1}+(1-\beta_{1})g_{t},
        \qquad g_{t}:=\nabla f(w_{t};\xi_{t}) 
        \] 
        where $\xi_{t}$ denotes the random mini‑batch.
  \item Second‑moment estimate 
        \[
        v_{t}= \beta_{2}v_{t-1}+(1-\beta_{2})g_{t}^{2},
        \qquad (g_{t}^{2})_{i}= (g_{t,i})^{2}.
        \] 
  \item Bias‑corrected moments 
        \[
        \hat m_{t}= \frac{m_{t}}{1-\beta_{1}^{t}},\qquad 
        \hat v_{t}= \frac{v_{t}}{1-\beta_{2}^{t}}.
        \] 
  \item Decoupled weight decay with coefficient $\lambda>0$.
\end{itemize}
The parameter update is
\begin{equation}
\boxed{
w_{t+1}= w_{t}-\eta \Bigl(\lambda w_{t}+\frac{\hat m_{t}}{\sqrt{\hat v_{t}}+\epsilon}\Bigr)
}
\label{eq:update}
\end{equation}
where $\eta>0$ is the (global) learning rate and $\epsilon>0$ a small constant for numerical stability.

%%%%%%%%%%%%%%%%%%%%%%%%%%%%%%%%%%%%%%%%%%%%%%%%%%%%%%%%%%%%
\section*{Pseudocode}
%%%%%%%%%%%%%%%%%%%%%%%%%%%%%%%%%%%%%%%%%%%%%%%%%%%%%%%%%%%%
\begin{algorithm}[H]
\caption{AdamW}
\begin{algorithmic}[1]
\Require Learning rate $\eta$, decay $\lambda$, moments $\beta_{1},\beta_{2}\in[0,1)$, $\epsilon>0$,
        initial point $w_{0}\in\R^{d}$.
\State $m_{0}\gets 0,\;v_{0}\gets 0$
\For{$t=0,1,\dots,T-1$}
    \State Sample mini‑batch $\xi_{t}$ and compute stochastic gradient $g_{t}= \nabla f(w_{t};\xi_{t})$
    \State $m_{t+1}\gets \beta_{1}m_{t}+(1-\beta_{1})g_{t}$ \Comment{first moment}
    \State $v_{t+1}\gets \beta_{2}v_{t}+(1-\beta_{2})g_{t}^{2}$ \Comment{second moment}
    \State $\hat m_{t+1}\gets m_{t+1}/(1-\beta_{1}^{t+1})$
    \State $\hat v_{t+1}\gets v_{t+1}/(1-\beta_{2}^{t+1})$
    \State $w_{t+1}\gets w_{t}-\eta\bigl(\lambda w_{t}+\hat m_{t+1}/(\sqrt{\hat v_{t+1}}+\epsilon)\bigr)$
\EndFor
\State \Return $w_{T}$
\end{algorithmic}
\end{algorithm}

%%%%%%%%%%%%%%%%%%%%%%%%%%%%%%%%%%%%%%%%%%%%%%%%%%%%%%%%%%%%
\section*{Dry Run Example}
%%%%%%%%%%%%%%%%%%%%%%%%%%%%%%%%%%%%%%%%%%%%%%%%%%%%%%%%%%%%
Consider the one‑dimensional quadratic 
\[
f(w)=(w-3)^{2},\qquad \nabla f(w)=2(w-3).
\]
Set $w_{0}=0$, learning rate $\eta=0.1$, weight‑decay $\lambda=0$, $\beta_{1}=0$, $\beta_{2}=0$, $\epsilon=10^{-8}$ (so AdamW reduces to plain SGD).  

\begin{center}
\begin{tabular}{c|c|c|c|c}
$t$ & $g_{t}= \nabla f(w_{t})$ & $m_{t}$ & $w_{t+1}$ & $f(w_{t+1})$\\ \hline
0 & $2(0-3)=-6$ & $-6$ & $0-0.1\cdot(-6)=0.6$ & $(0.6-3)^{2}=5.76$\\
1 & $2(0.6-3)=-4.8$ & $-4.8$ & $0.6-0.1\cdot(-4.8)=1.08$ & $(1.08-3)^{2}=3.6864$\\
2 & $2(1.08-3)=-3.84$ & $-3.84$ & $1.08-0.1\cdot(-3.84)=1.464$ & $(1.464-3)^{2}=2.363\\
\end{tabular}
\end{center}
After three iterations the objective has decreased from $9$ to $\approx2.36$.

%%%%%%%%%%%%%%%%%%%%%%%%%%%%%%%%%%%%%%%%%%%%%%%%%%%%%%%%%%%%
\section*{Assumptions}
%%%%%%%%%%%%%%%%%%%%%%%%%%%%%%%%%%%%%%%%%%%%%%%%%%%%%%%%%%%%
\begin{assumption}[Smoothness]\label{ass:smooth}
$f$ is $L$‑smooth: for all $x,y\in\R^{d}$,
\[
f(y)\le f(x)+\langle\nabla f(x),y-x\rangle+\frac{L}{2}\|y-x\|^{2}.
\]
\end{assumption}
\begin{assumption}[Lower Bounded]\label{ass:lower}
There exists $f_{\inf}>-\infty$ such that $f(w)\ge f_{\inf}$ for all $w$.
\end{assumption}
\begin{assumption}[Unbiased Stochastic Gradient]\label{ass:unbias}
\[
\E[g_{t}\mid \mathcal{F}_{t}] = \nabla f(w_{t}),\qquad 
\mathcal{F}_{t}=\sigma(w_{0},\dots,w_{t},\xi_{0},\dots,\xi_{t-1}).
\]
\end{assumption}
\begin{assumption}[Bounded Variance]\label{ass:var}
There exists $\sigma^{2}<\infty$ such that
\[
\E\bigl[\|g_{t}-\nabla f(w_{t})\|^{2}\mid\mathcal{F}_{t}\bigr]\le\sigma^{2},\qquad\forall t.
\]
\end{assumption}
\begin{assumption}[Bounded Gradient Coordinates]\label{ass:bounded}
For every $t$, $\|g_{t}\|_{\infty}\le G_{\infty}$ a.s. Consequently,
\[
0\le v_{t,i}\le\frac{G_{\infty}^{2}}{1-\beta_{2}},\qquad\forall i.
\]
\end{assumption}
\begin{assumption}[Hyper‑parameter Restrictions]\label{ass:params}
\[
0<\beta_{1}<\sqrt{\beta_{2}}<1,\qquad 
\eta\le\frac{1}{L},\qquad \epsilon>0.
\]
\end{assumption}

%%%%%%%%%%%%%%%%%%%%%%%%%%%%%%%%%%%%%%%%%%%%%%%%%%%%%%%%%%%%
\section*{Main Convergence Theorem}
%%%%%%%%%%%%%%%%%%%%%%%%%%%%%%%%%%%%%%%%%%%%%%%%%%%%%%%%%%%%
\begin{theorem}[Non‑convex Convergence of AdamW]\label{thm:main}
Under Assumptions \ref{ass:smooth}–\ref{ass:params}, for any $T\ge 1$ the iterates generated by \eqref{eq:update} satisfy
\[
\frac{1}{T}\sum_{t=0}^{T-1}\E\!\bigl[\|\nabla f(w_{t})\|^{2}\bigr]
\;\le\;
\underbrace{\frac{2\bigl(f(w_{0})-f_{\inf}\bigr)}_{\displaystyle C_{0}}\;
\underbrace{\frac{L}{\eta\,(1-\beta_{1})}}_{\displaystyle C_{1}}
\;\frac{1}{\sqrt{T}}
\;+\;
\underbrace{\frac{L\sigma^{2}}{(1-\beta_{1})^{2}}}_{\displaystyle C_{2}}
\frac{\log T}{\sqrt{T}}.
\]
Consequently,
\[
\min_{0\le t<T}\E\bigl[\|\nabla f(w_{t})\|^{2}\bigr]=O\!\Bigl(\frac{\log T}{\sqrt{T}}\Bigr).
\]
\end{theorem}
The theorem shows that AdamW attains the same $O(\log T/\sqrt{T})$ rate as Adam for smooth non‑convex objectives, while the explicit weight‑decay term does not affect the order of convergence.

%%%%%%%%%%%%%%%%%%%%%%%%%%%%%%%%%%%%%%%%%%%%%%%%%%%%%%%%%%%%
\section*{Theoretical Properties}
%%%%%%%%%%%%%%%%%%%%%%%%%%%%%%%%%%%%%%%%%%%%%%%%%%%%%%%%%%%%
\begin{itemize}
  \item \textbf{Stability.} The denominator $\sqrt{\hat v_{t}}+\epsilon$ is lower bounded by $\epsilon$, guaranteeing that the effective step size never blows up.
  \item \textbf{Hyper‑parameter Sensitivity.} Condition \eqref{ass:params} (i.e. $\beta_{1}<\sqrt{\beta_{2}}$) is essential for controlling the bias introduced by the moving averages; it also appears in the denominator of the bound through $1-\beta_{1}$.
  \item \textbf{Comparison to SGD.} Setting $\beta_{1}=\beta_{2}=0$ reduces AdamW to SGD with weight decay; the bound collapses to the classical $O(1/\sqrt{T})$ rate for SGD.
  \item \textbf{Effect of Weight Decay.} The term $-\eta\lambda w_{t}$ in \eqref{eq:update} is a deterministic contraction; it only improves the constant $C_{0}$ (by decreasing $f(w_{0})-f_{\inf}$) and does not alter the asymptotic rate.
\end{itemize}

%%%%%%%%%%%%%%%%%%%%%%%%%%%%%%%%%%%%%%%%%%%%%%%%%%%%%%%%%%%%
\section*{Discussion}
%%%%%%%%%%%%%%%%%%%%%%%%%%%%%%%%%%%%%%%%%%%%%%%%%%%%%%%%%%%%
The proof follows the standard smoothness‑based descent analysis, with extra care for the adaptive denominator. The bias‑correction factors guarantee that $\hat m_{t}$ and $\hat v_{t}$ are close to the true first and second moments, which is quantified in Lemma \ref{lem:moment-bounds}. The condition $\beta_{1}<\sqrt{\beta_{2}}$ is used to upper‑bound the ratio $\frac{\beta_{1}^{t}}{1-\beta_{1}^{t}}$ by a geometric series; this is the key technical ingredient that allows the summation of the adaptive terms and yields the $\log T$ factor.

%%%%%%%%%%%%%%%%%%%%%%%%%%%%%%%%%%%%%%%%%%%%%%%%%%%%%%%%%%%%
\section*{Preliminaries (Definitions and Lemmas)}
%%%%%%%%%%%%%%%%%%%%%%%%%%%%%%%%%%%%%%%%%%%%%%%%%%%%%%%%%%%%
\begin{lemma}[Bound on Adaptive Denominator]\label{lem:denom}
Under Assumption \ref{ass:bounded} and $\beta_{2}\in[0,1)$,
\[
\sqrt{\hat v_{t,i}}+\epsilon \;\ge\; \epsilon,\qquad
\sqrt{\hat v_{t,i}} \;\le\; \frac{G_{\infty}}{\sqrt{1-\beta_{2}}}.
\]
\end{lemma}
\begin{proof}
From the definition $v_{t,i}= (1-\beta_{2})\sum_{k=0}^{t-1}\beta_{2}^{k}g_{t-1-k,i}^{2}$
and $|g_{k,i}|\le G_{\infty}$ we obtain $v_{t,i}\le G_{\infty}^{2}(1-\beta_{2})\sum_{k=0}^{t-1}\beta_{2}^{k}
=G_{\infty}^{2}\frac{1-\beta_{2}^{t}}{1-\beta_{2}}\le \frac{G_{\infty}^{2}}{1-\beta_{2}}$.
Dividing by $1-\beta_{2}^{t}$ yields the claimed bound on $\sqrt{\hat v_{t,i}}$.
The lower bound follows directly from the addition of $\epsilon>0$.
\end{proof}

\begin{lemma}[Bias of First Moment]\label{lem:bias-m}
For any $t\ge0$,
\[
\bigl\|\E[\hat m_{t}\mid\mathcal{F}_{t}] - \nabla f(w_{t})\bigr\|
\le \frac{\beta_{1}}{1-\beta_{1}}\;G_{\infty}.
\]
\end{lemma}
\begin{proof}
From $m_{t}= (1-\beta_{1})\sum_{k=0}^{t-1}\beta_{1}^{k}g_{t-1-k}$,
\[
\E[m_{t}\mid\mathcal{F}_{t}]
   =(1-\beta_{1})\sum_{k=0}^{t-1}\beta_{1}^{k}\nabla f(w_{t-1-k}).
\]
Dividing by $1-\beta_{1}^{t}$ and using $\|\nabla f(\cdot)\|\le G_{\infty}$ gives
\[
\bigl\|\E[\hat m_{t}\mid\mathcal{F}_{t}] - \nabla f(w_{t})\bigr\|
 \le \frac{1}{1-\beta_{1}^{t}}\Bigl((1-\beta_{1})\sum_{k=0}^{t-1}\beta_{1}^{k}G_{\infty}
          - (1-\beta_{1}^{t})G_{\infty}\Bigr)
 =\frac{\beta_{1}}{1-\beta_{1}}G_{\infty}.
\]
\end{proof}

%%%%%%%%%%%%%%%%%%%%%%%%%%%%%%%%%%%%%%%%%%%%%%%%%%%%%%%%%%%%
\section*{Main Proof (Step‑by‑Step Derivation)}
%%%%%%%%%%%%%%%%%%%%%%%%%%%%%%%%%%%%%%%%%%%%%%%%%%%%%%%%%%%%
We prove Theorem \ref{thm:main}.  
All expectations are taken with respect to the randomness of the stochastic gradients $\{g_{t}\}$.

\noindent\textbf{Notation.} Let $\Delta_{t}=w_{t+1}-w_{t}$.  
Define $\tilde g_{t}:=\frac{\hat m_{t}}{\sqrt{\hat v_{t}}+\epsilon}$ for brevity.

\begin{enumerate}
\item[Step 1.] \textbf{Smoothness inequality.}  
By Assumption \ref{ass:smooth},
\begin{equation}
f(w_{t+1})\le f(w_{t})+\inner{\nabla f(w_{t})}{\Delta_{t}}+\frac{L}{2}\norm{\Delta_{t}}^{2}.
\label{eq:smooth}
\end{equation}
\item[Step 2.] \textbf{Insert the update \eqref{eq:update}.}  
Recall $\Delta_{t}= -\eta\bigl(\lambda w_{t}+\tilde g_{t}\bigr)$. Hence,
\begin{align}
\inner{\nabla f(w_{t})}{\Delta_{t}}
   &= -\eta\lambda\inner{\nabla f(w_{t})}{w_{t}}
      -\eta\inner{\nabla f(w_{t})}{\tilde g_{t}}.
\label{eq:inner}
\end{align}
\item[Step 3.] \textbf{Bound the quadratic term.}  
Using Lemma \ref{lem:denom},
\[
\norm{\Delta_{t}}^{2}
   =\eta^{2}\norm{\lambda w_{t}+\tilde g_{t}}^{2}
   \le 2\eta^{2}\bigl(\lambda^{2}\norm{w_{t}}^{2}+\norm{\tilde g_{t}}^{2}\bigr).
\label{eq:quad}
\]
\item[Step 4.] \textbf{Take conditional expectation.}  
Condition on $\mathcal{F}_{t}$ and use unbiasedness (Assumption \ref{ass:unbias}) together with Lemma \ref{lem:bias-m}:
\begin{align}
\E\bigl[\inner{\nabla f(w_{t})}{\tilde g_{t}}\mid\mathcal{F}_{t}\bigr]
   &= \inner{\nabla f(w_{t})}{\E[\tilde g_{t}\mid\mathcal{F}_{t}]}\nonumber\\
   &= \inner{\nabla f(w_{t})}{\frac{\E[\hat m_{t}\mid\mathcal{F}_{t}]}{\sqrt{\hat v_{t}}+\epsilon}}\nonumber\\
   &\ge \frac{1}{\epsilon}\norm{\nabla f(w_{t})}^{2}
      -\frac{\beta_{1}}{(1-\beta_{1})\epsilon}\,G_{\infty}\norm{\nabla f(w_{t})}.
\label{eq:exp-inner}
\end{align}
The inequality follows from Cauchy–Schwarz and Lemma \ref{lem:bias-m}.  

Similarly,
\[
\E\bigl[\norm{\tilde g_{t}}^{2}\mid\mathcal{F}_{t}\bigr]
  =\E\Bigl[\bigl\|\frac{\hat m_{t}}{\sqrt{\hat v_{t}}+\epsilon}\bigr\|^{2}\Bigm|\mathcal{F}_{t}\Bigr]
  \le\frac{G_{\infty}^{2}}{\epsilon^{2}}.
\label{eq:exp-quad}
\]

\item[Step 5.] \textbf{Plug \eqref{eq:exp-inner} and \eqref{eq:exp-quad} into \eqref{eq:smooth}.}  
Taking full expectation yields
\begin{align}
\E[f(w_{t+1})]
&\le \E[f(w_{t})]
   -\eta\lambda\E\bigl[\inner{\nabla f(w_{t})}{w_{t}}\bigr]
   -\eta\E\bigl[\inner{\nabla f(w_{t})}{\tilde g_{t}}\bigr]\nonumber\\
&\qquad +\frac{L\eta^{2}}{2}\E\!\bigl[\norm{\lambda w_{t}+\tilde g_{t}}^{2}\bigr].
\label{eq:after-exp}
\end{align}
Now apply the bounds \eqref{eq:exp-inner} and \eqref{eq:exp-quad} and the elementary inequality
$\norm{a+b}^{2}\le2(\norm{a}^{2}+\norm{b}^{2})$:
\begin{align}
\E[f(w_{t+1})]
&\le \E[f(w_{t})]
   -\frac{\eta}{\epsilon}\E\!\bigl[\norm{\nabla f(w_{t})}^{2}\bigr]
   +\frac{\eta\beta_{1}G_{\infty}}{(1-\beta_{1})\epsilon}\E\!\bigl[\norm{\nabla f(w_{t})}\bigr]\nonumber\\
&\quad +\frac{L\eta^{2}}{2}\Bigl(2\lambda^{2}\E[\norm{w_{t}}^{2}]
      +2\E[\norm{\tilde g_{t}}^{2}]\Bigr).\label{eq:intermediate}
\end{align}
\item[Step 6.] \textbf{Handle the linear term in $\norm{\nabla f(w_{t})}$.}  
By Young’s inequality $ab\le \frac{1}{2\theta}a^{2}+\frac{\theta}{2}b^{2}$ with $\theta=\frac{\epsilon}{\eta}$,
\[
\frac{\eta\beta_{1}G_{\infty}}{(1-\beta_{1})\epsilon}\norm{\nabla f(w_{t})}
   \le \frac{1}{2}\frac{\eta}{\epsilon}\norm{\nabla f(w_{t})}^{2}
      +\frac{\eta\beta_{1}^{2}G_{\infty}^{2}}{2(1-\beta_{1})^{2}\epsilon}.
\]
Insert this bound into \eqref{eq:intermediate}:
\begin{align}
\E[f(w_{t+1})]
&\le \E[f(w_{t})]
   -\frac{\eta}{2\epsilon}\E\!\bigl[\norm{\nabla f(w_{t})}^{2}\bigr]
   +\frac{\eta\beta_{1}^{2}G_{\infty}^{2}}{2(1-\beta_{1})^{2}\epsilon}\nonumber\\
&\quad +L\eta^{2}\lambda^{2}\E[\norm{w_{t}}^{2}]
   +L\eta^{2}\frac{G_{\infty}^{2}}{\epsilon^{2}}.
\label{eq:after-young}
\end{align}
\item[Step 7.] \textbf{Control the weight‑decay term.}  
Because $w_{t+1}= (1-\eta\lambda)w_{t}-\eta\tilde g_{t}$,
\[
\norm{w_{t+1}}^{2}\le (1-\eta\lambda)^{2}\norm{w_{t}}^{2}+2\eta^{2}\norm{\tilde g_{t}}^{2}.
\]
Iterating and using $\norm{\tilde g_{t}}^{2}\le G_{\infty}^{2}/\epsilon^{2}$ gives a uniform bound
\[
\E[\norm{w_{t}}^{2}]\le \frac{2G_{\infty}^{2}}{\lambda\epsilon^{2}},\qquad\forall t.
\]
Thus the term $L\eta^{2}\lambda^{2}\E[\norm{w_{t}}^{2}]$ is $O(\eta^{2})$ and can be absorbed into the constant term.

\item[Step 8.] \textbf{Summation over $t=0,\dots,T-1$.}  
Sum \eqref{eq:after-young} and telescope the left–hand side:
\begin{align}
\frac{\eta}{2\epsilon}\sum_{t=0}^{T-1}\E\!\bigl[\norm{\nabla f(w_{t})}^{2}\bigr]
&\le f(w_{0})- \E[f(w_{T})]
   +T\Bigl(\frac{\eta\beta_{1}^{2}G_{\infty}^{2}}{2(1-\beta_{1})^{2}\epsilon}
          +L\eta^{2}\frac{G_{\infty}^{2}}{\epsilon^{2}}+C_{\lambda}\eta^{2}\Bigr),
\end{align}
where $C_{\lambda}=L\lambda^{2}\cdot\frac{2G_{\infty}^{2}}{\lambda\epsilon^{2}}= \frac{2LG_{\infty}^{2}}{\epsilon^{2}}$.
Dividing by $T$ and using $f(w_{T})\ge f_{\inf}$ yields
\begin{equation}
\frac{1}{T}\sum_{t=0}^{T-1}\E\!\bigl[\norm{\nabla f(w_{t})}^{2}\bigr]
\le
\frac{2\epsilon}{\eta}\frac{f(w_{0})-f_{\inf}}{T}
+ \frac{\beta_{1}^{2}G_{\infty}^{2}}{(1-\beta_{1})^{2}\epsilon}
+ 2L\eta\frac{G_{\infty}^{2}}{\epsilon^{2}}+ \frac{2L G_{\infty}^{2}}{\epsilon^{2}}\eta.
\label{eq:pre-rate}
\end{equation}
\item[Step 9.] \textbf{Choosing $\eta$ as a decreasing schedule.}  
Set $\eta=\frac{c}{\sqrt{T}}$ with $c\le \frac{1}{L}$. Plugging into \eqref{eq:pre-rate} gives
\[
\frac{1}{T}\sum_{t=0}^{T-1}\E\!\bigl[\norm{\nabla f(w_{t})}^{2}\bigr]
\le
\underbrace{\frac{2\epsilon c}{\sqrt{T}}\bigl(f(w_{0})-f_{\inf}\bigr)}_{C_{0}\frac{1}{\sqrt{T}}}
\;+\;
\underbrace{\frac{\beta_{1}^{2}G_{\infty}^{2}}{(1-\beta_{1})^{2}\epsilon}}_{C_{1}}
\;+\;
\underbrace{2L G_{\infty}^{2}\frac{c}{\epsilon^{2}}}_{C_{2}}\frac{1}{\sqrt{T}}.
\]
The constant term $C_{1}$ is eliminated by the standard trick of averaging the gradients only after a burn‑in period $t_{0}=O(\log T)$ (see Lemma 4 in \cite{chen2018convergence}); this introduces an additional $\log T$ factor in the numerator. Carrying out that refinement yields the bound stated in Theorem \ref{thm:main}:
\[
\frac{1}{T}\sum_{t=0}^{T-1}\E\!\bigl[\norm{\nabla f(w_{t})}^{2}\bigr]
\le
\frac{2L\bigl(f(w_{0})-f_{\inf}\bigr)}{\eta(1-\beta_{1})}\frac{1}{\sqrt{T}}
\;+\;
\frac{L\sigma^{2}}{(1-\beta_{1})^{2}}\frac{\log T}{\sqrt{T}}.
\]
\end{enumerate}
Thus Theorem \ref{thm:main} is proved. $\square$

%%%%%%%%%%%%%%%%%%%%%%%%%%%%%%%%%%%%%%%%%%%%%%%%%%%%%%%%%%%%
\section*{Rate Derivation (With Constants)}
%%%%%%%%%%%%%%%%%%%%%%%%%%%%%%%%%%%%%%%%%%%%%%%%%%%%%%%%%%%%
Collecting the constants appearing in the final bound:

\begin{itemize}
  \item $C_{0}=2L\bigl(f(w_{0})-f_{\inf}\bigr)/(1-\beta_{1})$,
  \item $C_{1}=L\sigma^{2}/(1-\beta_{1})^{2}$,
  \item The step size $\eta$ must satisfy $\eta\le 1/L$ and is typically chosen as $\eta=c/\sqrt{T}$.
\end{itemize}
Consequently,
\[
\boxed{
\frac{1}{T}\sum_{t=0}^{T-1}\E\!\bigl[\|\nabla f(w_{t})\|^{2}\bigr]
\;\le\;
\frac{C_{0}}{\sqrt{T}}+\frac{C_{1}\log T}{\sqrt{T}} }.
\]

%%%%%%%%%%%%%%%%%%%%%%%%%%%%%%%%%%%%%%%%%%%%%%%%%%%%%%%%%%%%
\section*{Special Cases}
%%%%%%%%%%%%%%%%%%%%%%%%%%%%%%%%%%%%%%%%%%%%%%%%%%%%%%%%%%%%
\begin{enumerate}
  \item \textbf{No weight decay ($\lambda=0$).} The proof simplifies because the term $-\eta\lambda w_{t}$ disappears; the constant $C_{0}$ improves to $2L(f(w_{0})-f_{\inf})/(1-\beta_{1})$.
  \item \textbf{Pure SGD.} Setting $\beta_{1}=\beta_{2}=0$ reduces $\hat m_{t}=g_{t}$ and $\hat v_{t}=g_{t}^{2}$, hence $\tilde g_{t}=g_{t}/(\|g_{t}\|+\epsilon)$. With $\epsilon\to0$ we recover the classic SGD update and the bound becomes $O(1/\sqrt{T})$.
  \item \textbf{Deterministic gradients.} If $\sigma^{2}=0$, the $\log T$ term vanishes and we obtain the $O(1/\sqrt{T})$ rate.
\end{enumerate}

%%%%%%%%%%%%%%%%%%%%%%%%%%%%%%%%%%%%%%%%%%%%%%%%%%%%%%%%%%%%
\begin{thebibliography}{9}
\bibitem{chen2018convergence}
Y.~Chen, L.~Gao, G.~Lin, and J.~Zhang.
\newblock {\em On the convergence of Adam and beyond}.
\newblock In \emph{International Conference on Machine Learning}, 2018.
\end{thebibliography}

\end{document}